% Simple test document for CI/CD pipeline validation
\documentclass[11pt,a4paper]{article}

% Basic packages for XeLaTeX compatibility testing
\usepackage{fontspec}
\usepackage{geometry}
\usepackage{fancyhdr}

% Document setup
\geometry{a4paper, margin=2.5cm}
\setmainfont{Latin Modern Roman}

% Header and footer
\pagestyle{fancy}
\fancyhf{}
\fancyhead[L]{BRAT Annotation Guide}
\fancyhead[R]{CI/CD Pipeline Test}
\fancyfoot[C]{\thepage}

\title{CI/CD Pipeline Test Document}
\author{GitHub Actions}
\date{\today}

\begin{document}

\maketitle

\section{Pipeline Validation}

This document serves as a test for the GitHub Actions CI/CD pipeline. It validates:

\begin{itemize}
    \item XeLaTeX compilation with TeX Live 2024
    \item Font handling with fontspec package
    \item Basic document structure and formatting
    \item PDF generation and artifact upload
\end{itemize}

\subsection{Build Information}

If you see this document, the following pipeline components are working correctly:

\begin{enumerate}
    \item Repository checkout
    \item TeX Live 2024 installation
    \item XeLaTeX compilation
    \item PDF artifact generation
    \item Build logging and error handling
\end{enumerate}

\subsection{Template System Integration}

This test validates the basic LaTeX infrastructure required for the BRAT annotation guide template system.

\section{Next Steps}

Upon successful pipeline validation:
\begin{itemize}
    \item Main template compilation can proceed
    \item GitHub Pages dashboard will display build status
    \item Automated releases can be configured
    \item Full documentation pipeline is ready
\end{itemize}

\textbf{Pipeline Test Status:} \textcolor{green}{PASSED} ✓

\end{document}